\documentclass[10pt,a4paper]{amsart}
\usepackage{amsmath,amssymb,amsthm,graphicx}
% For commutative diagrams we use Xy-pic
\usepackage[all]{xy}

% We use 2cell for 2-commutative diagrams.
\xyoption{2cell}
% Theorem environments.
%
\theoremstyle{plain}
\newtheorem{theorem}[subsection]{Theorem}
\newtheorem{proposition}[subsection]{Proposition}
\newtheorem{lemma}[subsection]{Lemma}

\theoremstyle{definition}
\newtheorem{definition}[subsection]{Definition}
\newtheorem{example}[subsection]{Example}
\newtheorem{exercise}[subsection]{Exercise}
\newtheorem{situation}[subsection]{Situation}

\theoremstyle{remark}
\newtheorem{remark}[subsection]{Remark}
\newtheorem{remarks}[subsection]{Remarks}

\numberwithin{equation}{subsection}

% Macros
%
\def\lim{\mathop{\mathrm{lim}}\nolimits}
\def\colim{\mathop{\mathrm{colim}}\nolimits}
\def\Spec{\mathop{\mathrm{Spec}}}
\def\Hom{\mathop{\mathrm{Hom}}\nolimits}
\def\Ext{\mathop{\mathrm{Ext}}\nolimits}
\def\SheafHom{\mathop{\mathcal{H}\!\mathit{om}}\nolimits}
\def\SheafExt{\mathop{\mathcal{E}\!\mathit{xt}}\nolimits}
\def\Sch{\mathit{Sch}}
\def\Mor{\mathop{\mathrm{Mor}}\nolimits}
\def\Ob{\mathop{\mathrm{Ob}}\nolimits}
\def\Sh{\mathop{\mathit{Sh}}\nolimits}
\def\NL{\mathop{N\!L}\nolimits}
\def\CH{\mathop{\mathrm{CH}}\nolimits}
\def\proetale{{pro\text{-}\acute{e}tale}}
\def\etale{{\acute{e}tale}}
\def\QCoh{\mathit{QCoh}}
\def\Ker{\mathop{\mathrm{Ker}}}
\def\Im{\mathop{\mathrm{Im}}}
\def\Coker{\mathop{\mathrm{Coker}}}
\def\Coim{\mathop{\mathrm{Coim}}}

% Boxtimes
%
\DeclareMathSymbol{\boxtimes}{\mathbin}{AMSa}{"02}

%
% Macros for moduli stacks/spaces
%
\def\QCohstack{\mathcal{QC}\!\mathit{oh}}
\def\Cohstack{\mathcal{C}\!\mathit{oh}}
\def\Spacesstack{\mathcal{S}\!\mathit{paces}}
\def\Quotfunctor{\mathrm{Quot}}
\def\Hilbfunctor{\mathrm{Hilb}}
\def\Curvesstack{\mathcal{C}\!\mathit{urves}}
\def\Polarizedstack{\mathcal{P}\!\mathit{olarized}}
\def\Complexesstack{\mathcal{C}\!\mathit{omplexes}}
% \Pic is the operator that assigns to X its picard group, usage \Pic(X)
% \Picardstack_{X/B} denotes the Picard stack of X over B
% \Picardfunctor_{X/B} denotes the Picard functor of X over B
\def\Pic{\mathop{\mathrm{Pic}}\nolimits}
\def\Picardstack{\mathcal{P}\!\mathit{ic}}
\def\Picardfunctor{\mathrm{Pic}}
\def\Deformationcategory{\mathcal{D}\!\mathit{ef}}
\usepackage{cleveref}
\crefname{originallemma}{original lemma}{original lemmas}
\Crefname{originallemma}{Original lemma}{Original lemmas}
\makeatletter
\newlabel{coherent-lemma-quasi-coherent-affine-cohomology-zero}{{01XB}{source}{}{sourcepointer-1}{}}
\newlabel{coherent-lemma-quasi-coherent-affine-cohomology-zero@cref}{{[originallemma][0][]01XB}{[1][0][]source}{}{}{}}
\makeatother
\crefname{originallemma}{original lemma}{original lemmas}
\Crefname{originallemma}{Original lemma}{Original lemmas}
\makeatletter
\newlabel{descent-lemma-sheaf-condition-holds}{{03DT}{source}{}{sourcepointer-2}{}}
\newlabel{descent-lemma-sheaf-condition-holds@cref}{{[originallemma][0][]03DT}{[1][0][]source}{}{}{}}
\makeatother
\crefname{originalproposition}{original proposition}{original propositions}
\Crefname{originalproposition}{Original proposition}{Original propositions}
\makeatletter
\newlabel{descent-proposition-equivalence-quasi-coherent}{{03DX}{source}{}{sourcepointer-3}{}}
\newlabel{descent-proposition-equivalence-quasi-coherent@cref}{{[originalproposition][0][]03DX}{[1][0][]source}{}{}{}}
\makeatother
\crefname{originalproposition}{original proposition}{original propositions}
\Crefname{originalproposition}{Original proposition}{Original propositions}
\makeatletter
\newlabel{descent-proposition-same-cohomology-quasi-coherent}{{03DW}{source}{}{sourcepointer-4}{}}
\newlabel{descent-proposition-same-cohomology-quasi-coherent@cref}{{[originalproposition][0][]03DW}{[1][0][]source}{}{}{}}
\makeatother
\crefname{originalremark}{original remark}{original remarks}
\Crefname{originalremark}{Original remark}{Original remarks}
\makeatletter
\newlabel{descent-remark-change-topologies-ringed-sites}{{03FH}{source}{}{sourcepointer-5}{}}
\newlabel{descent-remark-change-topologies-ringed-sites@cref}{{[originalremark][0][]03FH}{[1][0][]source}{}{}{}}
\makeatother
\crefname{originalsection}{original section}{original sections}
\Crefname{originalsection}{Original section}{Original sections}
\makeatletter
\newlabel{descent-section-quasi-coherent-cohomology}{{0GN8}{source}{}{sourcepointer-6}{}}
\newlabel{descent-section-quasi-coherent-cohomology@cref}{{[originalsection][0][]0GN8}{[1][0][]source}{}{}{}}
\makeatother
\crefname{originalsection}{original section}{original sections}
\Crefname{originalsection}{Original section}{Original sections}
\makeatletter
\newlabel{descent-section-quasi-coherent-sheaves}{{03DR}{source}{}{sourcepointer-7}{}}
\newlabel{descent-section-quasi-coherent-sheaves@cref}{{[originalsection][0][]03DR}{[1][0][]source}{}{}{}}
\makeatother
\crefname{originalsection}{original section}{original sections}
\Crefname{originalsection}{Original section}{Original sections}
\makeatletter
\newlabel{descent-section-quasi-coherent-sheaves-bis}{{0GN9}{source}{}{sourcepointer-8}{}}
\newlabel{descent-section-quasi-coherent-sheaves-bis@cref}{{[originalsection][0][]0GN9}{[1][0][]source}{}{}{}}
\makeatother
\crefname{originalreference}{original reference}{original references}
\Crefname{originalreference}{Original reference}{Original references}
\makeatletter
\newlabel{proposition-quasi-coherent-sheaf-fpqc}{{03OG}{source}{}{sourcepointer-9}{}}
\newlabel{proposition-quasi-coherent-sheaf-fpqc@cref}{{[originalreference][0][]03OG}{[1][0][]source}{}{}{}}
\makeatother
\crefname{originalsection}{original section}{original sections}
\Crefname{originalsection}{Original section}{Original sections}
\makeatletter
\newlabel{section-cohomology-quasi-coherent}{{03OY}{source}{}{sourcepointer-10}{}}
\newlabel{section-cohomology-quasi-coherent@cref}{{[originalsection][0][]03OY}{[1][0][]source}{}{}{}}
\makeatother
\crefname{originalsection}{original section}{original sections}
\Crefname{originalsection}{Original section}{Original sections}
\makeatletter
\newlabel{section-quasi-coherent}{{01BD}{source}{}{sourcepointer-11}{}}
\newlabel{section-quasi-coherent@cref}{{[originalsection][0][]01BD}{[1][0][]source}{}{}{}}
\makeatother
\crefname{originallemma}{original lemma}{original lemmas}
\Crefname{originallemma}{Original lemma}{Original lemmas}
\makeatletter
\newlabel{sites-cohomology-lemma-higher-direct-images}{{072W}{source}{}{sourcepointer-12}{}}
\newlabel{sites-cohomology-lemma-higher-direct-images@cref}{{[originallemma][0][]072W}{[1][0][]source}{}{}{}}
\makeatother
\crefname{originalreference}{original reference}{original references}
\Crefname{originalreference}{Original reference}{Original references}
\makeatletter
\newlabel{theorem-quasi-coherent}{{03OJ}{source}{}{sourcepointer-13}{}}
\newlabel{theorem-quasi-coherent@cref}{{[originalreference][0][]03OJ}{[1][0][]source}{}{}{}}
\makeatother
\crefname{originalreference}{original reference}{original references}
\Crefname{originalreference}{Original reference}{Original references}
\makeatletter
\newlabel{theorem-zariski-fpqc-quasi-coherent}{{03P2}{source}{}{sourcepointer-14}{}}
\newlabel{theorem-zariski-fpqc-quasi-coherent@cref}{{[originalreference][0][]03P2}{[1][0][]source}{}{}{}}
\makeatother
\title{Quasi-coherent modules on the big etale and fppf sites of a scheme are the pullbacks of quasi-coherent modules on the scheme}
\author{The Stacks Project Authors}
\date{stacks-project:etale-cohomology.tex:local-snapshot-20261011; GNU FDL 1.2 or later}
\begin{document}
\maketitle
\noindent\textit{Editorial note.} All delivered mathematical text is the author's own, copied from the pinned Stacks Project chapter file etale-cohomology.tex; no statement or proof step was rewritten, reordered or removed.\par\medskip
\noindent\textit{Editorial note.} The counted proof is the author's complete argument: the identification of the presheaf F\textasciicircum{}a with the pullback of a quasi-coherent module on the scheme, the formal derivation of parts (2)-(6), and the computation of the higher direct images through an affine etale covering together with the comparison of fppf and Zariski cohomology.\par\medskip
\noindent\textit{Editorial note.} The label lemma-review-quasi-coherent names the counted statement, and the proof environment immediately following it in the source is delivered verbatim as its proof.\par\medskip
\noindent\textit{Editorial note.} The counted statement is carried by the author's own lemma environment in the chapter source.\par\medskip
\noindent\textit{Editorial note.} Pre-open gate: python operations/agent-authoring/tools/precard\_gate.py --source etale-cohomology.tex --label lemma-review-quasi-coherent --start 21926 --end 22025 --local-id 3565 returned RESULT=OK (L3565 free, occupancy FREE, public interval clear). Candidate taken from lane-s4/halfB-priority.json.\par\medskip
\noindent\textit{Editorial note.} The proof opens with the sentence 'We urge the reader to find their own proof of these results based on the material in Descent, Sections ...' — this is a reading suggestion, not a disclaimer: the proof environment that follows contains the argument (the identification of F\textasciicircum{}a with the pullback of a quasi-coherent module, the derivation of the formal parts and the computation of the higher direct images), so the statement is counted.\par\medskip
\noindent\textit{Editorial note.} The statement and proof are copied from the pinned chapter file; the cited inputs are kept as external pointers with their canonical Stacks tags.\par\medskip
\noindent\textit{Allowed standard prerequisites (editorial).}\begin{itemize}
\item Equivalence between quasi-coherent modules on the small Zariski site and on the big fpqc site. Descent, the proposition establishing the equivalence between quasi-coherent modules on a scheme and quasi-coherent modules on the big fpqc site, together with its variants for other topologies. Supplies the existence of the quasi-coherent module F\_0 on S whose pullback restricts to F, and the equivalence statements of parts (4) and (5).
\item Change of topology for ringed sites. Descent, the remark describing the pullback along the morphism of ringed sites comparing two topologies on the same category. Supplies the compatibility of the pullbacks under change of topology, which identifies F\textasciicircum{}a with F\_0\textasciicircum{}a and gives parts (2) and (3).
\item Descent of the sheaf condition along fpqc coverings. Descent, the lemma that the pullback of a quasi-coherent module satisfies the sheaf condition, together with the proposition on quasi-coherent sheaves for the fpqc topology. Supplies the sheaf property of F\textasciicircum{}a even for fpqc coverings, hence in particular for fppf and etale coverings.
\item Higher direct images as sheafified presheaves of cohomology, and vanishing of the cohomology of quasi-coherent sheaves on affine schemes. Cohomology on Sites, the lemma identifying the higher direct images with the sheaf associated to U maps to H\textasciicircum{}p(U, -); Descent, the comparison of fppf and Zariski cohomology for quasi-coherent modules; Cohomology of Schemes, the lemma that quasi-coherent sheaves on an affine scheme have vanishing higher cohomology. Supply the computation of R\textasciicircum{}p epsilon\_\{S,*\} and R\textasciicircum{}p a\_\{S,*\} on F\textasciicircum{}a through an affine etale covering of an arbitrary scheme over S.
\end{itemize}\medskip
\ifdefined\hypertarget\hypertarget{sourcepointer-1}{}\fi
\noindent\textit{Source reference.} Original lemma 01XB: sources/stacks-chow/tags.txt, source lines 2512--2512. Prerequisite pointer invoked at proof line(s) 22020; it names a standard fact of a sibling chapter that the counted argument uses, and every step of the counted proof built on it is retained verbatim in the delivered body.\par\medskip
\ifdefined\hypertarget\hypertarget{sourcepointer-2}{}\fi
\noindent\textit{Source reference.} Original lemma 03DT: sources/stacks-chow/tags.txt, source lines 4432--4432. Prerequisite pointer invoked at proof line(s) 21989; it names a standard fact of a sibling chapter that the counted argument uses, and every step of the counted proof built on it is retained verbatim in the delivered body.\par\medskip
\ifdefined\hypertarget\hypertarget{sourcepointer-3}{}\fi
\noindent\textit{Source reference.} Original proposition 03DX: sources/stacks-chow/tags.txt, source lines 4436--4436. Prerequisite pointer invoked at proof line(s) 21969; it names a standard fact of a sibling chapter that the counted argument uses, and every step of the counted proof built on it is retained verbatim in the delivered body.\par\medskip
\ifdefined\hypertarget\hypertarget{sourcepointer-4}{}\fi
\noindent\textit{Source reference.} Original proposition 03DW: sources/stacks-chow/tags.txt, source lines 4435--4435. Prerequisite pointer invoked at proof line(s) 22015; it names a standard fact of a sibling chapter that the counted argument uses, and every step of the counted proof built on it is retained verbatim in the delivered body.\par\medskip
\ifdefined\hypertarget\hypertarget{sourcepointer-5}{}\fi
\noindent\textit{Source reference.} Original remark 03FH: sources/stacks-chow/tags.txt, source lines 4492--4492. Prerequisite pointer invoked at proof line(s) 21985; it names a standard fact of a sibling chapter that the counted argument uses, and every step of the counted proof built on it is retained verbatim in the delivered body.\par\medskip
\ifdefined\hypertarget\hypertarget{sourcepointer-6}{}\fi
\noindent\textit{Source reference.} Original section 0GN8: sources/stacks-chow/tags.txt, source lines 20783--20783. Prerequisite pointer invoked at proof line(s) 21958; it names a standard fact of a sibling chapter that the counted argument uses, and every step of the counted proof built on it is retained verbatim in the delivered body.\par\medskip
\ifdefined\hypertarget\hypertarget{sourcepointer-7}{}\fi
\noindent\textit{Source reference.} Original section 03DR: sources/stacks-chow/tags.txt, source lines 4430--4430. Prerequisite pointer invoked at proof line(s) 21957; it names a standard fact of a sibling chapter that the counted argument uses, and every step of the counted proof built on it is retained verbatim in the delivered body.\par\medskip
\ifdefined\hypertarget\hypertarget{sourcepointer-8}{}\fi
\noindent\textit{Source reference.} Original section 0GN9: sources/stacks-chow/tags.txt, source lines 20784--20784. Prerequisite pointer invoked at proof line(s) 21959; it names a standard fact of a sibling chapter that the counted argument uses, and every step of the counted proof built on it is retained verbatim in the delivered body.\par\medskip
\ifdefined\hypertarget\hypertarget{sourcepointer-9}{}\fi
\noindent\textit{Source reference.} Original reference 03OG: sources/stacks-etale-cohomology/pinned-original-source.tex, source lines 1332--1332. Prerequisite pointer invoked at proof line(s) 21990; it names a standard fact of the same chapter that the counted argument uses, and every step of the counted proof built on it is retained verbatim in the delivered body.\par\medskip
\ifdefined\hypertarget\hypertarget{sourcepointer-10}{}\fi
\noindent\textit{Source reference.} Original section 03OY: sources/stacks-etale-cohomology/pinned-original-source.tex, source lines 2157--2157. Prerequisite pointer invoked at proof line(s) 21965; it names a standard fact of the same chapter that the counted argument uses, and every step of the counted proof built on it is retained verbatim in the delivered body.\par\medskip
\ifdefined\hypertarget\hypertarget{sourcepointer-11}{}\fi
\noindent\textit{Source reference.} Original section 01BD: sources/stacks-etale-cohomology/pinned-original-source.tex, source lines 1326--1326. Prerequisite pointer invoked at proof line(s) 21964; it names a standard fact of the same chapter that the counted argument uses, and every step of the counted proof built on it is retained verbatim in the delivered body.\par\medskip
\ifdefined\hypertarget\hypertarget{sourcepointer-12}{}\fi
\noindent\textit{Source reference.} Original lemma 072W: sources/stacks-chow/tags.txt, source lines 9040--9040. Prerequisite pointer invoked at proof line(s) 22007; it names a standard fact of a sibling chapter that the counted argument uses, and every step of the counted proof built on it is retained verbatim in the delivered body.\par\medskip
\ifdefined\hypertarget\hypertarget{sourcepointer-13}{}\fi
\noindent\textit{Source reference.} Original reference 03OJ: sources/stacks-etale-cohomology/pinned-original-source.tex, source lines 1416--1416. Prerequisite pointer invoked at proof line(s) 21996; it names a standard fact of the same chapter that the counted argument uses, and every step of the counted proof built on it is retained verbatim in the delivered body.\par\medskip
\ifdefined\hypertarget\hypertarget{sourcepointer-14}{}\fi
\noindent\textit{Source reference.} Original reference 03P2: sources/stacks-etale-cohomology/pinned-original-source.tex, source lines 2280--2280. Prerequisite pointer invoked at proof line(s) 22017; it names a standard fact of the same chapter that the counted argument uses, and every step of the counted proof built on it is retained verbatim in the delivered body.\par\medskip
\begin{lemma}
\label{lemma-review-quasi-coherent}
Let $S$ be a scheme. Let $\mathcal{F}$ be a quasi-coherent
$\mathcal{O}_S$-module on $S_\etale$.
\begin{enumerate}
\item The rule
$$
\mathcal{F}^a : (\Sch/S)_\etale \longrightarrow \textit{Ab},\quad
(f : T \to S) \longmapsto \Gamma(T, f_{small}^*\mathcal{F})
$$
satisfies the sheaf condition for fppf and a fortiori \'etale coverings,
\item $\mathcal{F}^a = \pi_S^*\mathcal{F}$ on $(\Sch/S)_\etale$,
\item $\mathcal{F}^a = a_S^*\mathcal{F}$ on $(\Sch/S)_{fppf}$,
\item the rule $\mathcal{F} \mapsto \mathcal{F}^a$ defines
an equivalence between quasi-coherent $\mathcal{O}_S$-modules
and quasi-coherent modules on
$((\Sch/S)_\etale, \mathcal{O})$,
\item the rule $\mathcal{F} \mapsto \mathcal{F}^a$ defines
an equivalence between quasi-coherent $\mathcal{O}_S$-modules
and quasi-coherent modules on
$((\Sch/S)_{fppf}, \mathcal{O})$,
\item we have $\epsilon_{S, *}a_S^*\mathcal{F} = \pi_S^*\mathcal{F}$
and $a_{S, *}a_S^*\mathcal{F} = \mathcal{F}$,
\item we have $R^i\epsilon_{S, *}(a_S^*\mathcal{F}) = 0$
and $R^ia_{S, *}(a_S^*\mathcal{F}) = 0$ for $i > 0$.
\end{enumerate}
\end{lemma}

\begin{proof}
We urge the reader to find their own proof of these results
based on the material in
Descent, Sections \ref{descent-section-quasi-coherent-sheaves},
\ref{descent-section-quasi-coherent-cohomology}, and
\ref{descent-section-quasi-coherent-sheaves-bis}.

\medskip\noindent
We first explain why the notation in this lemma is consistent with our
earlier use of the notation $\mathcal{F}^a$ in
Sections \ref{section-quasi-coherent} and
\ref{section-cohomology-quasi-coherent}
and in
Descent, Section \ref{descent-section-quasi-coherent-sheaves}.
Namely, we know by
Descent, Proposition \ref{descent-proposition-equivalence-quasi-coherent}
that there exists a quasi-coherent module
$\mathcal{F}_0$ on the scheme $S$ (in other words on the small
Zariski site) such that $\mathcal{F}$ is the restriction of the
rule
$$
\mathcal{F}_0^a : (\Sch/S)_\etale \longrightarrow \textit{Ab},\quad
(f : T \to S) \longmapsto \Gamma(T, f^*\mathcal{F})
$$
to the subcategory $S_\etale \subset (\Sch/S)_\etale$
where here $f^*$ denotes usual pullback of sheaves of modules on schemes.
Since $\mathcal{F}_0^a$ is pullback by the morphism of ringed
sites
$$
((\Sch/S)_\etale, \mathcal{O}) \longrightarrow (S_{Zar}, \mathcal{O}_{S_{Zar}})
$$
by Descent, Remark \ref{descent-remark-change-topologies-ringed-sites}
it follows immediately (from composition of pullbacks) that
$\mathcal{F}^a = \mathcal{F}_0^a$. This proves the sheaf property
even for fpqc coverings by
Descent, Lemma \ref{descent-lemma-sheaf-condition-holds} (see also
Proposition \ref{proposition-quasi-coherent-sheaf-fpqc}).
Then (2) and (3) follow
again by Descent, Remark \ref{descent-remark-change-topologies-ringed-sites}
and (4) and (5) follow from
Descent, Proposition \ref{descent-proposition-equivalence-quasi-coherent}
(see also the meta result
Theorem \ref{theorem-quasi-coherent}).

\medskip\noindent
Part (6) is immediate from the description of the sheaf
$\mathcal{F}^a = \pi_S^*\mathcal{F} = a_S^*\mathcal{F}$.

\medskip\noindent
For any abelian $\mathcal{H}$ on $(\Sch/S)_{fppf}$ the
higher direct image $R^p\epsilon_{S, *}\mathcal{H}$ is the sheaf
associated to the presheaf $U \mapsto H^p_{fppf}(U, \mathcal{H})$
on $(\Sch/S)_\etale$. See
Cohomology on Sites, Lemma \ref{sites-cohomology-lemma-higher-direct-images}.
Hence to prove
$R^p\epsilon_{S, *}a_S^*\mathcal{F} = R^p\epsilon_{S, *}\mathcal{F}^a = 0$
for $p > 0$ it suffices to show that any scheme $U$ over $S$
has an \'etale covering $\{U_i \to U\}_{i \in I}$ such that
$H^p_{fppf}(U_i, \mathcal{F}^a) = 0$ for $p > 0$.
If we take an open covering by affines, then the required
vanishing follows from comparison with usual cohomology
(Descent, Proposition \ref{descent-proposition-same-cohomology-quasi-coherent}
or
Theorem \ref{theorem-zariski-fpqc-quasi-coherent})
and the vanishing of cohomology of quasi-coherent sheaves
on affine schemes afforded by Cohomology of Schemes, Lemma
\ref{coherent-lemma-quasi-coherent-affine-cohomology-zero}.

\medskip\noindent
To show that $R^pa_{S, *}a_S^{-1}\mathcal{F} = R^pa_{S, *}\mathcal{F}^a = 0$
for $p > 0$ we argue in exactly the same manner. This finishes the proof.
\end{proof}

\end{document}
